% ---
% titre: Dossier d'exercices - Critères de divisibilité
% theme: NO
% chapitre: Nombres naturels et décimaux
% sous_chapitre: Critères de divisibilité
% niveau: [10e]
% type: EX
% tags: [c-criteres-divisibilite, f-individuel, f-maison, p-entrainement, p-evaluation-formative, p-institutionnalisation]
% difficulte: 3
% source: latex
% ---

\header{Les critères de divisibilité}

\vfill
{\bfseries\color{themeColor}Exercice 1.} Justifie en utilisant les critères de divisibilité.

945 est divisible par 5 car \sol{\hrulefill}{il finit par 5.}\\[6pt]
212 est divisible par 4 car \sol{\hrulefill}{12 est divisible par 4.}\\[6pt]
2\,168 est divisible par 8 car \sol{\hrulefill}{168 est divisible par 8.}

\vspace{30pt}
{\bfseries\color{themeColor}Exercice 2.} Complète le tableau suivant en notant une croix aux endroits qui conviennent.

\begin{center}
\def\arraystretch{2}
\begin{tabularx}{\linewidth}{|r*{7}{|>{\centering\arraybackslash}X}|}
\hline
\rowcolor{themeColor!20}\textbf{Divisible par \ldots} & \textbf{10} & \textbf{2} & \textbf{4} & \textbf{5} & \textbf{8} & \textbf{100} & \textbf{25}\\
\hline
1\,326 & & \textcolor{\colSolution} X & & & & & \\
\hline
37\,575 & & & & \textcolor{\colSolution} X & & & \textcolor{\colSolution} X \\
\hline
702\,500 & \textcolor{\colSolution} X & \textcolor{\colSolution} X & \textcolor{\colSolution} X & \textcolor{\colSolution} X & & \textcolor{\colSolution} X & \textcolor{\colSolution} X \\
\hline
432\,522 & & \textcolor{\colSolution} X & & & & & \\
\hline
4\,636 & & \textcolor{\colSolution} X & \textcolor{\colSolution} X & & \textcolor{\colSolution} X & & \\
\hline
\end{tabularx}
\def\arraystretch{1}
\end{center}

\vspace{30pt}
{\bfseries\color{themeColor}Exercice 3.} Écris tous les nombres divisibles par \og 4 \fg{} compris entre 400 et 440.

\sol{\bl{\linewidth}}{400, 404, 408, 412, 416, 420, 424, 428, 432, 436, 440}

\vspace{30pt}
{\bfseries\color{themeColor}Exercice 4.} Réponds par vrai (V) ou faux (F).

\def\arraystretch{1.5}
\begin{tabularx}{\linewidth}{lXl}
a. & Si un nombre est divisible par 4 alors il est divisible par 2. & \sol{\bl{1.2cm}}{V}\\
b. & Si un nombre est divisible par 2 et 5 alors il est divisible par 10. & \sol{\bl{1.2cm}}{V}\\
c. & Tous les nombres qui se terminent par 3 sont divisibles par 3. & \sol{\bl{1.2cm}}{F}\\
d. & Tout multiple de 10 est divisible par 2. & \sol{\bl{1.2cm}}{V}\\
e. & Si un nombre est divisible par 2 ou 4 alors il est divisible par 8. & \sol{\bl{1.2cm}}{F}\\
\end{tabularx}
\def\arraystretch{1}

\newpage
{\bfseries\color{themeColor}Exercice 5.} Par quel(s) chiffre(s) peux-tu remplacer l'espace vide pour que la phrase soit correcte ?

\def\arraystretch{1.5}
\begin{tabularx}{\linewidth}{lXl}
a. &75\_4  est divisible par 4. & \sol{\bl{10cm}}{0, 2, 4, 6 ou 8}\\
b. &3\,679\_  est divisible par 5 et par 2. & \sol{\bl{10cm}}{0}\\
c. &95\_2  est divisible par 8. & \sol{\bl{10cm}}{1, 5 ou 9}\\
d. &63\_0  est divisible par 10. & \sol{\bl{10cm}}{0, 1, 2, 3, 4, 5, 6, 7, 8, ou 9}\\
e. &972\_  est divisible par 25. &\sol{\bl{10cm}}{5}\\
\end{tabularx}
\def\arraystretch{1}

\vspace{30pt}
{\bfseries\color{themeColor}Exercice 6.} Voici une liste de nombres : recopie-les dans les cases s'ils sont divisibles par \ldots

\vspace{5pt}
\begin{center}
\itshape 7 \quad 80 \quad 35 \quad 200 \quad 375 \quad 600 \quad 472 \quad 108 \quad 420
\end{center}
\vspace{5pt}

\def\arraystretch{2}
\begin{tabularx}{\linewidth}{|l*{7}{|>{\centering\arraybackslash}X}|}
\hline
Par 2 &\textcolor{\colSolution}{80}&\textcolor{\colSolution}{200}&\textcolor{\colSolution}{600}&\textcolor{\colSolution}{472}&\textcolor{\colSolution}{108}&\textcolor{\colSolution}{420}& \\
\hline
Par 5 &\textcolor{\colSolution}{80}&\textcolor{\colSolution}{35}&\textcolor{\colSolution}{200}&\textcolor{\colSolution}{375}&\textcolor{\colSolution}{600}&\textcolor{\colSolution}{420}& \\
\hline
Par 10 &\textcolor{\colSolution}{80}&\textcolor{\colSolution}{200}&\textcolor{\colSolution}{600}&\textcolor{\colSolution}{420}&&& \\
\hline
Par 4 &\textcolor{\colSolution}{80}&\textcolor{\colSolution}{200}&\textcolor{\colSolution}{600}&\textcolor{\colSolution}{472}&\textcolor{\colSolution}{108}&\textcolor{\colSolution}{420}& \\
\hline
Par 8 &\textcolor{\colSolution}{80}&\textcolor{\colSolution}{200}&\textcolor{\colSolution}{600}&\textcolor{\colSolution}{472}&&& \\
\hline
Par 25 &\textcolor{\colSolution}{200}&\textcolor{\colSolution}{375}&\textcolor{\colSolution}{600}&&&& \\
\hline
\end{tabularx}
\def\arraystretch{1}

\vspace{30pt}
{\bfseries\color{themeColor}Exercice 7.} Relie chaque nombre à la bonne étiquette.

\vfill
\def\arraystretch{2}
\begin{tabularx}{\linewidth}{lrXlll}
a.& 5\,678&&\sol{}{b.}&1.&est divisible par 5 mais pas par 4.\\
b.& 7\,395&&\sol{}{c.}&2.&est divisible par 5 et par 4.\\
c.& 8\,760&&\sol{}{a.}&3.&n'est divisible ni par 5 ni par 4.\\
d.& 8\,976&&\sol{}{d.}&4.&est divisible par 4 mais pas par 5.\\
\end{tabularx}
\def\arraystretch{1}

\newpage
{\bfseries\color{themeColor}Exercice 8.} Relie chaque nombre au nombre qui le divise.

\def\arraystretch{2}
\begin{tabularx}{\linewidth}{lrXlll}
a.& 19\,432&&\sol{}{a.}&1.&est divisible par 4.\\
b.& 5\,046&&\sol{}{d.}&2.&est divisible par 10.\\
c.& 7\,7275&&\sol{}{c.}&3.&est divisible par 25.\\
d.& 3\,870&&\sol{}{b.}&4.&est divisible par 6.\\
\end{tabularx}
\def\arraystretch{1}

\vspace{30pt}
{\bfseries\color{themeColor}Exercice 9.} Relie chaque nombre au nombre qu'il divise.

\def\arraystretch{2}
\begin{tabularx}{\linewidth}{lrXlll}
a.& 5&&\sol{}{d.}&1.&divise 7\,896.\\
b.& 7&&\sol{}{b.}&2.&divise 7\,007.\\
c.& 9&&\sol{}{c.}&3.&divise 3\,681.\\
d.& 8&&\sol{}{a.}&4.&divise 4\,285.\\
\end{tabularx}
\def\arraystretch{1}

\vspace{30pt}
\noindent{\bfseries\color{themeColor}Exercice 10.} Souligne les nombres divisibles par 8.

\begin{center}
\sol{81\,936}{81\,936} \qquad 47\,284 \qquad 26\,774 \qquad \sol{70\,800}{70\,800} \qquad 83\,972 \qquad \sol{58\,576}{58\,576} \qquad \sol{13\,000}{13\,000}
\end{center}

\vspace{30pt}
{\bfseries\color{themeColor}Exercice 11.} Soustrais le plus petit nombre possible pour obtenir un nombre divisible par\ldots

\def\x{4cm}
\def\arraystretch{1.5}
\begin{tabularx}{\linewidth}{Xrcc}
Pour être divisible par 10 : & 827 &$-$ &\sol{\bl{\x}}{7} $=$ \sol{\bl{\x}}{820}\\
Pour être divisible par 8 : & 827 &$-$ &\sol{\bl{\x}}{3} $=$ \sol{\bl{\x}}{824}\\
Pour être divisible par 5 : & 827 &$-$ &\sol{\bl{\x}}{2} $=$ \sol{\bl{\x}}{825}\\
Pour être divisible par 4 : & 1\,138 &$-$ &\sol{\bl{\x}}{2} $=$ \sol{\bl{\x}}{1\,136}\\
Pour être divisible par 2 : & 1\,138 &$-$ &\sol{\bl{\x}}{0} $=$ \sol{\bl{\x}}{1\,138}\\
Pour être divisible par 8 : & 1\,138 &$-$ &\sol{\bl{\x}}{2} $=$ \sol{\bl{\x}}{1\,136}\\
\end{tabularx}
\def\arraystretch{1}

\newpage
{\bfseries\color{themeColor}Exercice 12.} Souligne les nombres qui conviennent.

\def\arraystretch{2}
\begin{tabularx}{\linewidth}{Xccccccc}
Divisible par 2: & 589 & \sol{13\,190}{13\,190} & \sol{7\,000}{7\,000} & \sol{841\,602}{841\,602} & \sol{19\,822}{19\,822} & 85\,369 & \sol{6\,378}{6\,378}\\
Divisible par 4: & \sol{3\,628}{3\,628} & \sol{60\,000}{60\,000} & \sol{4\,120}{4\,120} & \sol{876}{876} & \sol{9\,544}{9\,544} & 70\,289 & \sol{7\,652}{7\,652}\\
Divisible par 5: & \sol{6\,350}{6\,350} & \sol{745}{745} & \sol{14\,000}{14\,000} & 873 & \sol{74\,165}{74\,165} & 16\,276 & \sol{29\,310}{29\,310}\\
Divisible par 10: & \sol{6\,840}{6\,840} & 89\,735 & \sol{152\,000}{152\,000} & 89 & \sol{66\,560}{66\,560} & 4\,964 & \sol{54\,860}{54\,860}\\
Divisible par 25: & 645 & \sol{80\,000}{80\,000} & 3\,770 & \sol{450}{450} & \sol{41\,225}{41\,225} & 61\,077 & \sol{76\,800}{76\,800}\\
Divisible par 2 et 4: & 3\,822 & \sol{620}{620} & \sol{14\,000}{14\,000} & \sol{78\,636}{78\,636} & \sol{23\,612}{23\,612} & 10\,778 & \sol{75\,340}{75\,340}\\
Divisible par 2 et 5: & 598 & \sol{91\,740}{91\,740} & \sol{70\,000}{70\,000} & 24\,876 & \sol{45\,070}{45\,070} & 95\,659 & \sol{37\,790}{37\,790}\\
Divisible par 5 et 25: & 872 & \sol{14\,750}{14\,750} & 345 & \sol{75}{75} & \sol{15\,525}{15\,525} & 54\,854 & \sol{65\,150}{65\,150}\\
\end{tabularx}
\def\arraystretch{1}

\vspace{30pt}
{\bfseries\color{themeColor}Exercice 13.} Donne le nombre supérieur le plus proche de 1\,479 qui sera divisible par \ldots

\def\arraystretch{2}
\begin{tabularx}{\linewidth}{lcXlc}
2 & \sol{\bl{7cm}}{1\,480} && 5 & \sol{\bl{7cm}}{1\,480}\\
4 & \sol{\bl{7cm}}{1\,480} && 10 & \sol{\bl{7cm}}{1\,480}\\
8 & \sol{\bl{7cm}}{1\,480}&& 25 & \sol{\bl{7cm}}{1\,500}\\
\end{tabularx}
\def\arraystretch{1}

\vspace{30pt}
\noindent{\bfseries\color{themeColor}Exercice 14.} Avec les chiffres 4, 5 et 8, écris \ldots

\def\arraystretch{2}
\begin{tabularx}{\linewidth}{Xc}
Le plus grand nombre divisible par 2 : & \sol{\bl{10cm}}{854}\\
Le plus grand nombre divisible par 4 : & \sol{\bl{10cm}}{584}\\
Le plus grand nombre divisible par 8 : & \sol{\bl{10cm}}{548}\\
Le plus grand nombre divisible par 5 : & \sol{\bl{10cm}}{845}\\
\end{tabularx}
\def\arraystretch{1}